% ========================================================
% CONCLUSIONES
% ========================================================
\section{Conclusiones}

\subsection{Conclusiones del grupo}

La literatura reciente acelera K-means por dos vías. Una reparte el mismo trabajo entre más hilos o
máquinas; la otra hace menos trabajo mediante poda, otra representación del cálculo o
aproximación. Solo la primera conserva por construcción el resultado del algoritmo secuencial, y
es la que implementaremos en PC2.

En todas las escalas revisadas, el costo de sincronizar condiciona el diseño. Entre máquinas se
mide en rondas y en volumen comunicado; en una sola máquina, en contención sobre el estado
compartido. Nuestro diseño lo acota con acumuladores privados y una única barrera por iteración.

Las cifras de aceleración de los artículos van de 1,3 a 123,8, y buena parte de esa distancia se
explica por la línea base elegida. Nuestra evaluación fijará como referencia la versión secuencial
del mismo algoritmo, con la misma inicialización y la media recortada de varias ejecuciones.

El dataset ya está listo para el modelo: cada regla de limpieza está justificada, la verifican dos
implementaciones independientes y todo el proceso se reproduce con un comando.

\subsection{Conclusiones individuales}

\subsubsection*{\nombreRody{} (Tema 3)}
\pendiente{análisis y conclusiones propias sobre poda, localidad y aproximación, y su relación con
el caso}

\subsubsection*{\nombreDayana{} (Tema 2)}
\pendiente{análisis y conclusiones propias sobre sincronización distribuida, y su relación con el
caso}

\subsubsection*{\nombreTercero{} (Tema 1)}
\pendiente{análisis y conclusiones propias sobre paralelismo en memoria compartida, y su relación
con el caso}
